\section*{अध्याय \ref{ch:b2:eruptive-styles} --- ज्वालामुखी उद्गार}

\begin{solution}{exo:b2:eruptive-styles:1}
(क) $\qty{1.3}{km^3} = \qty{1.3e9}{m^3}$, $10^9$ और $10^{10}$ के बीच: VEI~5।
(ख) $\qty{2.9e8}{m^3}$: VEI~4। (ग) $\qty{2e6}{m^3}$, $10^6$ और $10^7$ के बीच:
VEI~2। (घ) VEI~6 \qty{e10}{m^3} से शुरू होता है, इसलिए स्थूल \omterm{def:b2:eruptive-styles:tephra}{टेफ़्रा} कम-से-कम
\qty{10}{km^3} होना चाहिए; ``\qty{5}{km^3} से अधिक'' एक निचली सीमा है, विरोधाभास नहीं।
\end{solution}

\begin{solution}{exo:b2:eruptive-styles:2}
(क) राख (\qty{2}{mm} से महीन)। (ख) एक लैपिलस (\qtyrange{2}{64}{mm})।
(ग) एक ब्लॉक: \qty{64}{mm} से बड़ा, कोणीय, फेंके जाते समय ठोस। (घ) एक
\omterm{def:b2:eruptive-styles:tephra}{ज्वालामुखीय बम}, एक तकुआ बम: \qty{64}{mm} से बड़ा और उड़ान में गलित रहते हुए
आकार लिया हुआ।
\end{solution}

\begin{solution}{exo:b2:eruptive-styles:3}
(क) $0.304\times(10^6)^{0.241} = 0.304\times27.9 = \qty{8.5}{km}$;
$0.304\times(10^8)^{0.241} = 0.304\times84.7 = \qty{25.8}{km}$।
(ख) $Q = (35/0.304)^{1/0.241} = 115.1^{4.149} = \qty{3.6e8}{kg/s}$, और $I =
8.55 + 3 = 11.6$। (ग) $Q \propto H^{4.15}$, इसलिए \qty{20}{\percent} अधिक ऊँचाई
$Q$ को $1.2^{4.15} = 2.1$ के गुणक से अधिक आँकती है।
\end{solution}

\begin{solution}{exo:b2:eruptive-styles:4}
$V_{\mathrm{DRE}} = 0.29\times1000/2500 = \qty{0.116}{km^3}$; $m =
\num{0.29e9}\times1000 = \qty{2.9e11}{kg}$; $M = \log_{10}(\num{2.9e11}) - 7 = 11.46
- 7 = 4.46$।
\end{solution}

\begin{solution}{exo:b2:eruptive-styles:5}
$k = \ln(40/10)/(15 - 5) = 1.386/10 = \qty{0.139}{km^{-1}}$; $T_0 = 40\,e^{0.139\times5}
= 40\times2.00 = \qty{80}{cm}$; $\ln2/k = \qty{5.0}{km}$। $T_0 =
\qty{8.0e-4}{km}$ के साथ: $V = 2\times\num{8.0e-4}/0.139^2 = \qty{0.083}{km^3} =
\qty{8.3e7}{m^3}$, VEI~3।
\end{solution}

\begin{solution}{exo:b2:eruptive-styles:6}
(क) $\int_0^x T_0e^{-ky}2y\,dy = 2T_0\big[-\tfrac{y}{k}e^{-ky} - \tfrac{1}{k^2}e^{-ky}\big]_0^x
= \tfrac{2T_0}{k^2}[1 - (1 + kx)e^{-kx}]$। (ख) $u = \ln(2(1 + u))$ की पुनरावृत्ति
$u = 1$ से: $1.386$, $1.564$, $1.635$, $1.662$, $1.672$, \dots, जो
$u = 1.678$ पर अभिसरित होती है, जिसके लिए $(1 + u)e^{-u} = 2.678\times0.1867 = 0.500$। (ग) $x_{1/2} =
1.678/0.139 = \qty{12.1}{km}$, सबसे बाहरी समस्थूलता-रेखा ($\sqrt A =
\qty{15}{km}$) के भीतर; फिर भी उस रेखा से परे का आयतन $(1 + 2.09)e^{-2.09} =
\qty{38}{\percent}$ है कुल का। निक्षेप का लगभग दो-पाँचवाँ भाग वहाँ है जहाँ
कुछ भी मापा नहीं गया: आयतन बहिर्वेशन पर टिका है।
\end{solution}

\begin{solution}{exo:b2:eruptive-styles:7}
(क) $p_s = (5.0/0.444)^2 = \qty{127}{MPa}$, जो $127/0.025 = \qty{5.1}{km}$ पर है।
(ख) $x = (5.0 - 0.444\sqrt{10})/100 = (5.0 - 1.404)/100 = 0.0360$;
$V_g = 0.0360\times8.314\times1123/(0.018015\times10^7) = \qty{1.86e-3}{m^3}$;
$V_m = 0.964/2300 = \qty{4.19e-4}{m^3}$; $\phi = 1.86/(1.86 + 0.419) = 0.82$।
(ग) बेसाल्टी गलित रायोलाइटी गलित से हज़ारों से अरबों गुना कम श्यान होता है
(\cref{ch:b2:magma-emplacement}): उसके बुलबुले उठते हैं,
मिलते हैं और निकल जाते हैं, और गैस एक पारगम्य झाग से होकर या छिद्र पर
निकल जाती है, इसलिए \omterm{def:b2:magma-emplacement:magma}{मैग्मा} अपनी गैस रोके रखने के बजाय खुले तंत्र में विगैसित होता है; उसका
विखंडन, जब होता है, छिद्र के निकट एक फ़व्वारे तक सीमित रहता है। धीमा
आरोहण इसके लिए अधिक समय देता है।
\end{solution}

\begin{solution}{exo:b2:eruptive-styles:8}
(क) $u_z(0) = 0.75\times10^7/(\pi\times(5000)^2) = \qty{0.095}{m} = \qty{95}{mm}$।
(ख) $z_0 = 3.8/0.766 = \qty{5.0}{km}$। (ग) अधिकतम का $(1 + 1)^{-3/2} = 0.35$
भाग।
\end{solution}

\begin{solution}{exo:b2:eruptive-styles:9}
अधिकतम परास \ang{45} पर है: $R = v^2/g = 100^2/9.81 = \qty{1.0}{km}$, और
$200^2/9.81 = \qty{4.1}{km}$। वायु-प्रतिरोध उड़ान को छोटा करता है, इसलिए ये
ऊपरी सीमाएँ हैं, बड़े खंडों के लिए सत्य के निकट और छोटों के लिए बहुत अधिक लंबी।
चूँकि $R = v^2/g$, संकट-क्षेत्र प्रक्षेपण-चाल के वर्ग के अनुसार बढ़ता है:
विस्फोट की चाल दोगुनी करने से उसकी त्रिज्या चौगुनी हो जाती है।
\end{solution}

\begin{solution}{exo:b2:eruptive-styles:10}
(क) $T = T_0e^{-k_1x}$ के साथ, जब $x < X$, और उसके परे $T = T_Xe^{-k_2(x - X)}$, जहाँ $T_X =
T_0e^{-k_1X}$ (सांतत्य), $\int_0^X ye^{-ky}dy = \frac{1}{k^2} - e^{-kX}\big(\frac{X}{k} +
\frac{1}{k^2}\big)$ और $\int_X^\infty ye^{-ky}dy = e^{-kX}\big(\frac{X}{k} +
\frac{1}{k^2}\big)$ का प्रयोग कीजिए:
निकटस्थ भाग $2T_0\big[\frac{1}{k_1^2} - e^{-k_1X}\frac{k_1X + 1}{k_1^2}\big]$ देता है और
दूरस्थ भाग $2T_Xe^{k_2X}\int_X^\infty ye^{-k_2y}dy = 2T_X\frac{k_2X + 1}{k_2^2}$; $T_X =
T_0e^{-k_1X}$ के साथ उनका योग विधि-बॉक्स का सूत्र है।
(ख) $T_0 = \qty{0.01285}{km}$: $2T_0/k_1^2 = 0.0257/0.1332 = \qty{0.193}{km^3}$;
$e^{-0.365\times8.56} = 0.0440$; $(0.138\times8.56 + 1)/0.138^2 = 114.5$ और
$(0.365\times8.56 + 1)/0.365^2 = 30.96$; संशोधन $0.0257\times0.0440\times83.6
= \qty{0.094}{km^3}$ है और $V = \qty{0.287}{km^3}$, प्रकाशित \qty{0.29}{km^3}।
$X$ से परे $2T_X(X/k_2 + 1/k_2^2) = 2\times\num{5.65e-4}\times114.5 =
\qty{0.129}{km^3}$ पड़ता है, अर्थात \qty{45}{\percent}।
\end{solution}

\begin{solution}{exo:b2:eruptive-styles:11}
(क) $100^{1/4} = 3.16$; आनुभविक नियम $100^{0.241} = 3.03$ और
$100^{0.226} = 2.83$ देते हैं। (ख) $H \propto N^{-3/4}$: अधिक प्रबलता से स्तरित
वायुमंडल स्तंभ को नीचे रोक देता है; $N$ के \qty{20}{\percent} बड़ा होने पर,
$1.2^{-0.75} = 0.87$, यानी \qty{13}{\percent} कम ऊँचाई। (ग) दुर्बल प्लूम
हवा से झुक जाते हैं, जो उन्हें नीचा करता है, जबकि नम वायु उनमें गुप्त ऊष्मा छोड़ती है,
जो उन्हें ऊँचा करती है; ये प्रभाव छोटे और बड़े उद्गारों पर भिन्न रूप से भारी पड़ते हैं
और सरल नियम में नहीं हैं। और दरें केवल कई के
गुणकों के भीतर ज्ञात हैं, क्षैतिज अक्ष के चर में एक त्रुटि, जो
बिठाई गई ढाल को शून्य की ओर झुका देती है।
\end{solution}

\begin{solution}{exo:b2:eruptive-styles:12}
(क) 1900 से पहले $132/400\times100 = 33$ प्रति शताब्दी; उसके बाद $95/125\times100 = 76$ प्रति
शताब्दी। (ख) $33/76 = \qty{43}{\percent}$। (ग) $21/56 = \qty{38}{\percent}$:
वही कम-अभिलेखन, दो से तीन के गुणक से। (घ) बड़े उद्गार
अधिक दूर से दिखते हैं और अधिक मोटे, चौड़े निक्षेप छोड़ते हैं जो अपरदन से अधिक समय तक
बचे रहते हैं; 1700 में किसी सुदूर द्वीप पर VEI~4 का उद्गार बिना दर्ज हुए निकल सकता था,
VEI~6 का उद्गार शायद ही।
\end{solution}

\begin{solution}{pb:b2:eruptive-styles:1}
\textbf{1.} $\ln T = 5.298$, $4.605$, $3.871$, $2.639$, $1.386$। क्रमिक समस्थूलता-रेखाओं के बीच
ढालें $-0.0693$, $-0.0367$, $-0.0308$ और $-0.0313$ प्रति km हैं: 10 और \qty{20}{km}
के बीच एक खड़ा खंड, 40 से \qty{120}{km} तक एक धीमा;
विराम 20 और \qty{40}{km} के बीच है।

\textbf{2.} $k_1 = \ln(200/100)/10 = \qty{0.0693}{km^{-1}}$; $T_0 = 200\,e^{0.693} =
\qty{400}{cm}$।

\textbf{3.} $167e^{-1.244} = \qty{48.1}{cm}$, $167e^{-2.488} = \qty{13.9}{cm}$,
$167e^{-3.732} = \qty{4.00}{cm}$: 48, 14 और 4.0 के \qty{1}{\percent} के भीतर।

\textbf{4.} $400e^{-0.0693X} = 167e^{-0.0311X}$ से $X = \ln(400/167)/(0.0693 -
0.0311) = 0.873/0.0382 = \qty{22.9}{km}$ मिलता है।

\textbf{5.} $\ln2/0.0693 = \qty{10.0}{km}$ और $\ln2/0.0311 = \qty{22.3}{km}$।

\textbf{6.} $T_0 = \qty{0.0040}{km}$: $V = 2\times0.0040/0.0693^2 =
\qty{1.67}{km^3}$।

\textbf{7.} $e^{-0.0693\times22.9} = 0.205$; $(0.0311\times22.9 + 1)/0.0311^2 = 1770$;
$(0.0693\times22.9 + 1)/0.0693^2 = 539$; संशोधन $2\times0.0040\times0.205\times1231 =
\qty{2.02}{km^3}$; $V = 1.67 + 2.02 = \qty{3.68}{km^3}$, जिसमें संशोधन
\qty{55}{\percent} है।

\textbf{8.} $3.68/1.67 = 2.2$। महीन दूरस्थ राख पतली होती है, आसानी से अपरदित,
पुनः-संचलित या समुद्र में गिरी हुई; किसी पुराने निक्षेप में वह प्रायः खो जाती है या कभी
मानचित्रित नहीं होती, और तब आयतन दो या अधिक के गुणक से कम आँका जाता है।

\textbf{9.} $\qty{3.68e9}{m^3}$ $10^9$ और $10^{10}$ के बीच है: VEI~5।

\textbf{10.} $V_{\mathrm{DRE}} = 3.68\times900/2400 = \qty{1.38}{km^3}$।

\textbf{11.} $m = \num{3.68e9}\times900 = \qty{3.31e12}{kg}$।

\textbf{12.} $M = \log_{10}(\num{3.31e12}) - 7 = 12.52 - 7 = 5.52$।

\textbf{13.} $3.68/1.3 = 2.8$ गुना।

\textbf{14.} \qty{25}{km} से ऊँचा स्तंभ, और प्लिनियन शैली।

\textbf{15.} $Q = \num{3.31e12}/36000 = \qty{9.2e7}{kg/s}$; $I = 7.96 + 3 = 11.0$।

\textbf{16.} $0.304\times(\num{9.2e7})^{0.241} = 0.304\times83.0 = \qty{25.2}{km}$।

\textbf{17.} $0.345\times(\num{9.2e7})^{0.226} = 0.345\times63.1 = \qty{21.8}{km}$।

\textbf{18.} मापा गया \qty{31}{km} पहले पूर्वानुमान से \qty{23}{\percent} ऊपर है
और दूसरे से \qty{42}{\percent} ऊपर, नियमों के
\qtyrange{53}{60}{\percent} प्रकीर्णन के भीतर: संगत।

\textbf{19.} $Q = (31/0.304)^{1/0.241} = 102.0^{4.149} = \qty{2.2e8}{kg/s}$, माध्य दर का 2.3
गुना। खंड सबसे ऊँचे स्तंभ को दर्ज करते हैं, उद्गार के शिखर पर,
जबकि प्रश्न~15 पूरे चरण पर औसत लेता है; और द्रव्यमान का कुछ भाग
निक्षेप से गायब हो सकता है (दूरस्थ राख, \omterm{def:b2:eruptive-styles:pdc}{पायरोक्लास्टिक प्रवाह})।

\textbf{20.} हाँ: 22 से \qty{31}{km}, \qty{16}{km} पर क्षोभसीमा से काफ़ी ऊपर।
समतापमंडल में पहुँची गैस पूरी पृथ्वी के चारों ओर फैलती है; मुख्य गैस
सल्फ़र डाइऑक्साइड है, और 1991 में पिनातूबो द्वारा वहाँ पहुँचाई गई दो करोड़ टन ने
पृथ्वी को दो वर्षों तक लगभग आधा डिग्री ठंडा किया।

\textbf{21.} $\lambda_3 = 399/75 = 5.32$, $\lambda_4 = 95/125 = 0.760$, $\lambda_5 = 22/225 =
0.0978$, $\lambda_6 = 8/525 = 0.0152$ प्रति वर्ष; लघुगणक $0.726$, $-0.119$, $-1.010$,
$-1.817$।

\textbf{22.} रेखा $10^{0.73} = 5.37$, $10^{-0.12} = 0.759$, $10^{-0.97} =
0.107$ और $10^{-1.82} = 0.0151$ देती है: गणनाओं के क्रमशः \qty{1}{\percent}, \qty{0}{\percent},
\qty{10}{\percent} और \qty{1}{\percent} के भीतर।

\textbf{23.} बड़े उद्गारों का अभिलेख समय में और पीछे तक पूर्ण रहता है,
क्योंकि वे अधिक दूर से देखे जाते हैं और उनके निक्षेप अधिक समय तक बचे रहते हैं; 1500 में
शुरू होने वाली खिड़की VEI~3 के उद्गारों को बुरी तरह कम गिनती, जबकि 1950 में
शुरू होने वाली खिड़की में VEI~6 का लगभग कोई उद्गार न होता।

\textbf{24.} $\lambda_5 = \num{0.107}$ प्रति वर्ष, यानी पुनरावृत्ति काल $1/0.107 =
\num{9.3}\text{ वर्ष}$; दस वर्षों में $1 - e^{-1.07} = 1 - 0.343 = \qty{66}{\percent}$।

\textbf{25.} $\lambda_7 = 10^{3.28 - 5.95} = 10^{-2.67} = 0.00214$ प्रति वर्ष, हर
470 वर्षों में एक; अभिलेख $2/1025 = 0.00195$ प्रति वर्ष देता है: \qty{10}{\percent} के भीतर
सहमति, उस बिंदु के लिए जो फ़िट में प्रयुक्त नहीं हुआ।

\textbf{26.} $\lambda_6 = 10^{3.28 - 5.10} = 10^{-1.82} = 0.0151$ प्रति वर्ष; $\lambda t = 1.51$
और $1 - e^{-1.51} = 1 - 0.220 = 0.780$: अगले सौ वर्षों में पृथ्वी पर कहीं
VEI~6 या अधिक के कम-से-कम एक उद्गार की प्रायिकता
\textbf{\qty[text-series-to-math=true]{78}{\percent}} है। यह एक स्थिर दर पर,
एक-दूसरे से स्वतंत्र उद्गारों (एक प्वासों प्रक्रम) पर, और अपनी खिड़कियों में
पूर्ण सूची पर टिकी है।
\end{solution}
